\section{Introduction}
Mobile and embedded platforms such as smartphones, portable media players, and tablets have repeatedly been shown to suffer poor user experience from memory pressure, fragmentation, and wasteful reclamation~\cite{TCE_Lim2013_VMP, TC_10492614}. In a managed runtime, the region of memory that stores application objects is the \textit{heap}; it is grown and shrunk automatically by the allocator, and object reclamation is delegated to a garbage collector rather than to explicit deallocation. Modern Android systems address heap pressure by combining the Android Runtime (ART) with compacting generational garbage collection, compressed-memory mechanisms such as zRAM, and kernel-level controls originally explored in embedded and virtualized devices~\cite{TCE_Hwang2012_CompressedSwap, TC_9140304, TC_10851391}. Despite these mechanisms, heap footprint still grows in ways the runtime cannot anticipate, because native (C/C++) libraries reached through the Java Native Interface (JNI)---for example, media, graphics, and inference frameworks---allocate outside the managed heap and are therefore invisible to ART's accounting. This unmanaged growth is a major cause of latency spikes, forced termination of background processes, and performance degradation under load.

Modern mobile devices thus combine a managed Java heap under the runtime with an unmanaged native heap accessed through C/C++ libraries. This dual-heap architecture reflects long-standing trade-offs between performance and control in general-purpose and embedded computing~\cite{TCE_Shin2010_FTL,TC_9252863}. While the runtime reclaims Java objects through concurrent garbage collection, native allocations remain largely opaque to it, producing asymmetric visibility and delayed response under sustained workloads such as real-time media processing, interactive graphics, and on-device inference~\cite{TC_9513533}. As pressure builds, kernel-level reclaim and swap are triggered, frequently causing excessive reclamation overhead and unpredictable application behavior~\cite{TC_10851391}.

Earlier Android generations already shifted heap-management responsibility from application developers to the managed runtime, mirroring trends across mobile and embedded consumer platforms. Runtime-centric regulation, however, is fundamentally \emph{reactive}: ART begins collecting only after a generation threshold is crossed, and the kernel begins reclaiming only after page-level pressure is detected, so both mechanisms respond to overshoot rather than preventing it. Although Linux cgroup controllers and memory limits provide isolation, they lack semantic awareness of application-level heap dynamics, especially when multiple libraries share memory pages. This mismatch between kernel-level reclamation and runtime accounting often results in premature process termination or repeated garbage-collection cycles~\cite{TCE_Kim2018_QoS, TC_11313797}. The problem is intensified by on-device intelligence, where memory behavior interacts tightly with inference pipelines, on-device learning, and adaptive workloads~\cite{TCE_Protogeros2025_FL,TC_10564838}; existing approaches, however, optimize mainly for computational or communication efficiency and do not integrate heap growth and reclamation into a single OS-level control loop.

\smallskip
\noindent\textbf{Positioning.} We group prior memory-management work into three families and state up front how HMS differs from each. \emph{Runtime-level} techniques adapt garbage collection to allocation patterns~\cite{TC_10045641,TCE_Yan2019_FileAwareGC} but reason only about managed objects and cannot see native growth. \emph{Kernel-level} techniques reclaim and compress pages~\cite{TCE_Hwang2012_CompressedSwap,TC_9140304,TC_9513533} but act on physical pages without object- or heap-level semantics. \emph{Hardware-assisted} mechanisms such as ARM Memory Tagging detect memory-safety errors but do not regulate heap growth. A few recent efforts coordinate across tiers for edge and cloud inference~\cite{TC_10320372,TCE_Wang2025_CloudEdge}, yet they prioritize network efficiency and latency rather than heap stability. HMS differs by making the runtime's logical view of heap growth and the kernel's physical reclaim actuators \emph{observe and act through one shared control loop}, so that reclamation is both semantically informed and temporally aligned. Section~\ref{sec_related} contrasts these families in detail (Table~\ref{tab:related_compare}).

To close the visibility and timing gaps identified above, we propose an OS-level \textbf{Heap Memory Safeguard (HMS)} that coordinates heap control across the runtime and kernel layers. We first analyze empirical, cross-layer characteristics of heap behavior on commercial devices to expose the persistent visibility and timing gaps between runtime garbage collection and kernel-level reclamation~\cite{TC_10492614,TC_10045641, TC_9140304}. Building on these insights, HMS bridges the kernel's physical memory view and the runtime's logical heap by monitoring allocation trends and coordinating targeted reclamation before pressure escalates. The design draws on prior work in kernel-aware compaction and workload-sensitive memory management~\cite{TCE_Lim2013_VMP,TC_9513533} and extends it to Android-class platforms through explicit cross-layer coordination. This coordination problem rarely arises in general-purpose systems: mobile and CE devices operate under strict memory budgets, often lack swap-backed virtual memory, and must sustain predictable interactive latency while continuously running tightly coupled managed and native workloads.

\smallskip
\noindent\textbf{Contributions.} This paper makes the following contributions:
\begin{enumerate}
  \item \textbf{Empirical characterization.} Using commercial devices, we quantify a \emph{visibility gap} (ART is blind to native heap growth until a page fault) and a \emph{temporal gap} (runtime and kernel act reactively and at different times), and we show that these gaps, not raw memory scarcity alone, drive interactive stalls (Section~\ref{sec_motivation}).
  \item \textbf{A fused utilization signal with an explicit rationale.} We introduce the heap-utilization index $H_u$, which combines resident footprint with allocation velocity, and we derive why this combination separates transient bursts from sustained growth better than footprint-only, rate-only, or pressure-only signals (Section~\ref{subsec_rationale}).
  \item \textbf{A cross-layer system, HMS.} We design and implement HMS---a Runtime Observer, a Kernel Enforcer, and a memcg/tagging feedback interface---that requires no application changes and reuses standard interfaces (Section~\ref{sec_hms}).
  \item \textbf{Analysis that isolates the source of the benefit.} Beyond aggregate speedups, we analyze \emph{why} HMS helps: we examine the reclaim--GC interaction, quantify control overhead and misprediction behavior, and separate the contribution of cross-layer coordination from that of merely reclaiming more aggressively (Sections~\ref{sec_eval}--\ref{sec_case}).
\end{enumerate}
